% El playbook de usuario, en español. El aspecto está en sheriff-playbook.sty,
% compartido con la versión en inglés (sheriff-playbook-en.tex): un cambio en
% lo que el usuario escribe o lee va a las dos en el mismo cambio.
% scripts/playbook.sh compila ambas.
\documentclass[11pt,a4paper]{article}
\usepackage[spanish]{sheriff-playbook}

\begin{document}

\makecover{Sheriff}{Estándares de código de Kaizten\\en Claude Code, Codex y Antigravity}{Instalación, uso diario, integración continua y resolución de problemas.}

\section*{Cómo usar este playbook}

Este playbook es para quien escribe código con Claude Code, Codex o Antigravity y quiere
que ese código cumpla los estándares de Kaizten desde la primera línea. No hace
falta saber cómo están hechas las herramientas: aquí se cuenta qué instalar, qué
pedir y qué esperar en cada momento.

Las secciones 1 a 3 se leen una vez, al empezar. A partir de la 4, el playbook se
organiza en \emph{jugadas}: cada una es una situación concreta, con cuándo
usarla, qué pedir, qué ocurre y cómo saber que ha terminado. La última página es
una tarjeta de referencia para tener a mano.

\begin{tip}
\textbf{Si solo tienes cinco minutos.}
\begin{enumerate}[topsep=2pt,leftmargin=1.3em]
  \item Ejecuta la línea de instalación; no hace falta cuenta de GitHub. Si
        falta algo, como Java 17 o Docker, el instalador te lo dice
        (sección~\ref{sec:instalacion}). En Linux y macOS:
\end{enumerate}
\begin{command}[basicstyle=\ttfamily\scriptsize]{bash}
curl -fsSL https://github.com/kaizten/sheriff-mcp-server/releases/latest/download/install.sh \
  | bash
\end{command}
\noindent\hspace*{1.3em}En Windows, en PowerShell:
\begin{command}[basicstyle=\ttfamily\footnotesize]{PowerShell}
$u = 'https://github.com/kaizten/sheriff-mcp-server/releases/latest/download';
& ([scriptblock]::Create((irm "$u/install.ps1")))
\end{command}
\begin{enumerate}[resume,leftmargin=1.3em]
  \item Abre Claude Code, Codex o Antigravity en tu proyecto y pide
        \ask{revisa los estándares de código}.
\end{enumerate}
\end{tip}

\vspace{0.2\baselineskip}
{\sffamily\bfseries\Large\color{pbbrand}Contenido}\par
\vspace{-0.7\baselineskip}
{\color{pbrule}\rule{\linewidth}{0.8pt}}
\makeatletter\@starttoc{toc}\makeatother

\clearpage
% ---------------------------------------------------------------------------
\section{Qué es Sheriff y qué hace por ti}
\label{sec:que-es}

\textbf{Sheriff} es el analizador de estándares de código de Kaizten. Comprueba,
por ejemplo, que cada método tenga su JavaDoc, que los imports estén ordenados,
que cada \code{if} lleve llaves o que el código respete la arquitectura del
proyecto (hexagonal, DDD). Sus reglas van compiladas en una imagen de Docker,
\code{kaizten/sheriff}: no son una guía de estilo que interpretar, sino
comprobaciones que pasan o fallan.

Estas herramientas llevan Sheriff hasta el sitio donde escribes código. Con
ellas, tu asistente:

\begin{itemize}
  \item comprueba el componente \textbf{antes} de tocarlo y limpia los errores
        que ya había;
  \item consulta las reglas antes de escribir, para no introducir errores nuevos;
  \item aplica primero las \textbf{reparaciones automáticas} de Sheriff, que no
        gastan tokens, y corrige a mano solo lo que queda;
  \item no da el trabajo por terminado mientras queden errores o fallen los tests
        del proyecto.
\end{itemize}

\begin{center}
\begin{tikzpicture}[node distance=4mm and 7mm]
  \node[pbstrong,minimum width=17mm,minimum height=13mm] (you) {Tú};
  \node[pbbox,right=of you,minimum width=29mm,minimum height=13mm] (assistant) {\textbf{Tu asistente}\\[1pt]\scriptsize Claude Code, Codex\\\scriptsize o Antigravity};
  \node[pbbox,right=15mm of assistant,minimum width=37mm] (server) {\textbf{Servidor MCP}\\[1pt]\scriptsize cinco herramientas};
  \node[pbbox,below=of server,minimum width=37mm] (hooks) {\textbf{Hooks}\\[1pt]\scriptsize antes de editar, al terminar};
  \node[pbbox,below=of hooks,minimum width=37mm] (agent) {\textbf{Agente de reparación}\\[1pt]\scriptsize \code{sheriff_autofix}, \code{--agent}};
  \node[pbbox,below=of agent,minimum width=37mm] (check) {\textbf{Comprobación en CI}\\[1pt]\scriptsize \code{--check}};
  \begin{scope}[on background layer]
    \node[draw=pbaccent,dashed,line width=0.9pt,fill=pbaccentsoft,rounded corners=5pt,inner sep=3.5mm,fit=(server)(hooks)(agent)(check)] (jar) {};
  \end{scope}
  \node[anchor=south,font=\sffamily\scriptsize\bfseries,text=pbaccent] at (jar.north) {sheriff-mcp.jar};
  \node[pbstrong,right=11mm of jar,minimum width=23mm,minimum height=15mm] (sheriff) {Sheriff\\[1pt]\scriptsize\color{pbaccentsoft}imagen Docker};
  \node[pbsoft,minimum width=29mm] (ci) at (assistant |- check) {CI, terminal\\Maven};
  \draw[pbarrow] (you) -- (assistant);
  \draw[pbarrow] (assistant) -- node[pblabel,above]{MCP} (server);
  \draw[pbarrow,draw=pbmuted] (hooks.west) -| node[pblabel,pos=0.3,below,align=center]{te paran si\\hay errores} (assistant.south);
  \draw[pbarrow,draw=pbmuted] (ci) -- (check);
  \draw[pbarrow] (jar.east) -- (sheriff.west);
\end{tikzpicture}
\end{center}

Todo vive en un único archivo, \code{sheriff-mcp.jar}, que se instala una vez
por máquina. Tu proyecto no necesita copiar nada ni configurarse: las
herramientas deducen qué analizar y con qué reglas a partir de la carpeta en la
que trabajas.

\subsection{Lo que cambia en tu día a día}

\begin{itemize}
  \item \textbf{Pides los cambios como siempre.} El asistente comprueba los
        estándares por su cuenta; no hace falta pedírselo.
  \item \textbf{Si el componente ya tenía errores, los arregla primero}, aunque
        no los hayas introducido tú. Es intencionado: el código que se entrega
        queda limpio entero. Si en algún caso no lo quieres, dilo en la
        petición: \ask{sin tocar los errores que ya había}.
  \item \textbf{Antes de terminar, comprueba otra vez y pasa los tests} del
        proyecto, y te dice qué errores estaban antes de empezar.
  \item \textbf{Si intenta terminar con errores} en lo que ha cambiado, un hook
        le devuelve el turno. No tienes que vigilarlo tú.
\end{itemize}

\subsection{Seis palabras que conviene conocer}

\begin{center}
\rowcolors{2}{pbcode}{white}
\begin{tabularx}{\linewidth}{@{\hspace{6pt}}K{33mm}L}
\headrow \thead{Palabra} & \thead{Qué significa aquí} \\
Componente & La carpeta que Sheriff analiza: un módulo del proyecto, el que
             tiene \code{src/}. En un proyecto de un solo módulo, el propio
             proyecto. \\
Perfil & El conjunto de reglas que se aplica: \code{JAVA},
         \code{JAVA_HEXAGONAL}, \code{TYPESCRIPT}\ldots{} Se deduce solo. \\
Error & Un incumplimiento de una regla, con su archivo, su línea y su código
        (por ejemplo, \code{SortedImport}). \\
Reparación automática & Un \emph{fixer} de Sheriff: corrige una regla sin
        ningún modelo de IA y sin gastar tokens. \\
Hook & Un comando que Claude Code o Codex ejecutan solos en un momento
       concreto: antes de editar un archivo o al terminar un turno. \\
Tarea & Cada análisis queda registrado con un identificador
        (\code{Task 4047f176.}), útil para consultar o informar de algo. \\
\end{tabularx}
\end{center}

% ---------------------------------------------------------------------------
\section{Instalación}
\label{sec:instalacion}

Una sola línea instala la última versión publicada, en Linux, macOS y Windows.
\textbf{La misma línea actualiza} cuando sale una nueva. El instalador comprueba
lo que necesita y, si falta algo, te dice qué hacer.

\subsection{Lo único que necesitas antes}
\label{sec:requisitos}

Nada: ni cuenta de GitHub ni ninguna herramienta suya. La línea de
instalación descarga el script de la página pública de releases del
repositorio, y el script descarga el resto de la misma. Java y Docker tampoco
hace falta tenerlos antes: el instalador los busca y te dice qué hacer.

\subsection{La línea de instalación}

\needspace{4\baselineskip}
En Linux y macOS:

\begin{command}[basicstyle=\ttfamily\scriptsize]{bash}
curl -fsSL https://github.com/kaizten/sheriff-mcp-server/releases/latest/download/install.sh | bash
\end{command}

Sin \code{curl}, \code{wget -qO-} con la misma dirección hace lo mismo. En
Windows, en PowerShell (Windows PowerShell o PowerShell 7):

\begin{command}[basicstyle=\ttfamily\footnotesize]{PowerShell}
$u = 'https://github.com/kaizten/sheriff-mcp-server/releases/latest/download';
& ([scriptblock]::Create((irm "$u/install.ps1")))
\end{command}

\subsection{Si prefieres leer el script antes}

La línea de instalación pasa el script directamente a \code{bash}. Si prefieres
descargarlo, leerlo y ejecutarlo después, hazlo así; \code{--release} le dice
que instale la release aunque lo ejecutes dentro de un repositorio:

\begin{command}[basicstyle=\ttfamily\scriptsize]{bash}
curl -fsSLO https://github.com/kaizten/sheriff-mcp-server/releases/latest/download/install.sh
bash install.sh --release
\end{command}

\begin{command}[basicstyle=\ttfamily\footnotesize]{PowerShell}
$u = 'https://github.com/kaizten/sheriff-mcp-server/releases/latest/download';
irm "$u/install.ps1" -OutFile install.ps1
.\install.ps1 -Release
\end{command}

Si PowerShell se niega a ejecutar el archivo, permite los scripts solo en esa
ventana con \path{Set-ExecutionPolicy -Scope Process -ExecutionPolicy Bypass}.
Los dos scripts también se pueden descargar desde la página de releases del
repositorio en GitHub.

\subsection{Lo que comprueba el instalador}

\begin{center}
\rowcolors{2}{pbcode}{white}
\begin{tabularx}{\linewidth}{@{\hspace{6pt}}K{33mm}L}
\headrow \thead{Comprueba} & \thead{Si falta} \\
Java 17 o posterior & Lo busca en \code{JAVA_HOME}, en el \code{PATH} y en
  las carpetas donde se instalan los JDK, aunque no esté en el \code{PATH}, y
  configura los asistentes con su ruta completa. Si no hay ninguno, en Windows
  instala Temurin 21 con \code{winget}; en Linux y macOS se detiene y te indica
  dónde descargarlo (\path{https://adoptium.net}). \\
Docker & Sigue sin él y te avisa de que no se podrá analizar nada hasta que lo
  tengas. Con \code{--install-docker} lo instala (en Linux, con el script de
  Docker o, en Arch, con \code{pacman}, y te añade al grupo \code{docker}; en
  macOS y Windows, Docker Desktop) y lo arranca si está parado. \\
Claude Code & Si está, registra el servidor MCP. Si no, te imprime el comando
  para registrarlo cuando lo instales. \\
Codex & Si está, lo configura (sección~\ref{sec:codex}). Si no, te imprime el
  comando para cuando lo instales. \\
Antigravity & Si está \code{agy}, lo configura (sección~\ref{sec:antigravity}).
  Si no, te imprime el comando para cuando lo instales. \\
El servidor MCP & Al final lo arranca una vez, como lo arrancará tu asistente,
  y comprueba que responde: \code{OK}, o \code{FAIL} con la causa. \\
\end{tabularx}
\end{center}

Los hooks de Claude Code los escribe aunque \code{claude} no esté en el
\code{PATH}, si existe la carpeta \path{~/.claude}: son solo un archivo de
ajustes, y la aplicación de escritorio y las extensiones de los editores
también lo leen. No toca ningún proyecto. Las reglas de Sheriff (el \emph{catálogo}) no
vienen en la descarga: las herramientas las extraen de la imagen la primera vez
que las necesitan, en unos 6 segundos.

\begin{warning}
Docker Desktop pide aceptar sus condiciones la primera vez que arranca, y el
instalador no las acepta por ti. En empresas de más de 250 personas o de más de
10 millones de dólares de facturación anual exigen una suscripción de pago.
\end{warning}

\subsection{Lo que verás}

Esta es la salida de una instalación real en una máquina con Claude Code,
Codex y Antigravity. Entre la tercera línea y la cuarta aparece también la descarga de la
imagen de Sheriff: unos 4\,GB la primera vez y, después, solo lo que cambie.

\begin{answer}[Salida del instalador]
Using Java 17 at /usr/bin/java
Downloading sheriff-mcp.jar from the latest release of kaizten/sheriff-mcp-server...
Installed sheriff-mcp 1.0.2 into /home/ana/.local/share/sheriff-agent
Sheriff's hooks are in /home/ana/.claude/settings.json, for every project Claude Code opens:
  before an edit     refuse new code on a component that already fails the standards
  when a turn ends   refuse to finish while what the turn changed has errors
  when a turn starts note what was already changed, so that it is not the turn's
sheriff_test, sheriff_fix, sheriff_guidelines and sheriff_task run without asking.
Where there is nothing Sheriff analyzes, they let everything through.
Undo with: java -jar sheriff-mcp.jar --uninstall-hooks --user
Added stdio MCP server sheriff with command: /usr/bin/java -jar /home/ana/.local/share/sheriff-agent/sheriff-mcp.jar to user config
File modified: /home/ana/.claude.json
Registered the MCP server as 'sheriff'. Restart open Claude Code sessions to pick it up.
Codex is set up for every project it opens:
  MCP server     'sheriff' in /home/ana/.codex/config.toml, with 600 s per call and sheriff_fix approved
  order of work  in /home/ana/.codex/AGENTS.md
Undo with: java -jar sheriff-mcp.jar --uninstall-codex
Sheriff's hooks are in /home/ana/.codex/hooks.json, for every project Codex opens:
  when a turn ends   refuse to finish while what the turn changed has errors
  when a turn starts note what was already changed, so that it is not the turn's
Codex runs them once you approve them: it asks the first time a session starts, or use /hooks.
Until then a turn can end with errors left: approve them in your first session.
Restart open Codex sessions to pick it up; the first one asks you to review the new hooks.
Antigravity is set up for every project agy opens:
  MCP server     'sheriff' in /home/ana/.gemini/config/mcp_config.json
  permissions    sheriff_test, sheriff_fix, sheriff_guidelines and sheriff_task run without asking,
                 in /home/ana/.gemini/antigravity-cli/settings.json
It has no Sheriff hooks: a turn there can end with errors left.
Undo with: java -jar sheriff-mcp.jar --uninstall-antigravity
Restart open agy sessions to pick it up.
The MCP server starts: OK, it answers as sheriff.

Sheriff is set up for Claude Code, Codex and Antigravity.
\end{answer}

Dice qué Java usará, qué versión ha instalado y dónde y, para cada asistente, qué hooks ha
activado, qué herramientas pueden ejecutarse sin pedir permiso y cómo deshacer
cada cosa; después, que el servidor arranca, y en su última línea, qué
asistentes ha configurado. Solo configura
los que la máquina ya tiene, y nunca instala ninguno: un asistente que no está
recibe una línea que dice cómo configurarlo cuando lo esté. Como pide,
\textbf{reinicia las sesiones de Claude Code y de Codex que tuvieras abiertas}.
Para comprobar que tus asistentes ven el servidor: en Claude Code,
\code{claude mcp list} debe mostrar
\code{sheriff:}~\ldots~\textcolor{pbtip}{\faCheck}~\code{Connected} (dentro de una
sesión, \code{/mcp}); en Codex, \code{codex mcp list} debe mostrar \code{sheriff}
como \code{enabled}, y en Antigravity, \code{agy mcp list}, igual.

\subsection{Opciones}

Las opciones van al final de la línea de instalación: en bash, detrás de
\code{bash -s --}; en PowerShell, detrás del bloque. Por ejemplo:

\begin{command}{bash}
curl -fsSL .../install.sh | bash -s -- --pull-always
\end{command}
\begin{command}{PowerShell}
& ([scriptblock]::Create((irm "$u/install.ps1"))) -PullAlways
\end{command}

\begin{center}
\rowcolors{2}{pbcode}{white}
\begin{tabularx}{\linewidth}{@{\hspace{6pt}}>{\ttfamily\small}p{31mm}>{\ttfamily\small}p{26mm}L}
\headrow \thead{bash} & \thead{PowerShell} & \thead{Qué hace} \\
-{}-install-docker & -InstallDocker & Instala Docker si falta y lo arranca si
  está parado. \\
 & -DockerWsl & Windows: usa el Docker instalado dentro de WSL, sin Docker
  Desktop. Después, reinicia Claude Code. \\
-{}-pull-always & -PullAlways & Descarga una imagen de Sheriff más nueva cada vez
  que arranca el servidor. Sin ella, solo se avisa una vez por sesión, para que
  no cambien las reglas a mitad de un trabajo. Hay que repetirla en cada
  actualización. \\
-{}-no-pull & -NoPull & No descarga la imagen ahora; el servidor la descarga
  cuando la necesite. \\
-{}-no-hooks & -NoHooks & Instala el servidor MCP sin los hooks. \\
-{}-no-codex & -NoCodex & No toca la configuración de Codex. \\
-{}-no-antigravity & -NoAntigravity & No toca la configuración de Antigravity. \\
-{}-fail-fast & -FailFast & Los hooks paran Sheriff en el primer error: van más
  rápido, pero solo informan de uno. \\
\end{tabularx}
\end{center}

% ---------------------------------------------------------------------------
\section{Tu primera sesión}
\label{sec:primera}

Este es un ejemplo real, de principio a fin, con un proyecto Maven muy pequeño
llamado \code{basket}. Las respuestas de Sheriff están copiadas tal cual,
recortadas solo donde aparece \code{[...]}.

\subsubsection{1. Abre el asistente en el proyecto}
Abre una terminal en la raíz del proyecto (la carpeta con el \code{pom.xml},
el \code{package.json} o el \code{src/}) y arranca \code{claude}, \code{codex} o
\code{agy}: las respuestas de Sheriff y los pasos son los mismos. Desde el
editor o la aplicación de escritorio vale igual, siempre que la carpeta abierta
sea la del proyecto.

\subsubsection{2. Pide la revisión}
Escribe \ask{revisa los estándares de código}. El asistente llama a
\tool{sheriff_test}, que responde con un recuento:

\begin{answer}
basket has 4 error(s) under JAVA.
Paths below are under basket/src/main/java/demo/
  Basket.java: 4

By rule: JavaJavaDocCommentInMethodMessage6 2, JavaBracesIfStatementsMessage2 1, SortedImport 1

Task 4047f176.
Next step: call the sheriff_fix tool with component 'basket'. It applies every repair Sheriff has for these errors and lists what it cannot repair.
\end{answer}

\begin{note}
\textbf{Cómo leer una respuesta.} La primera línea dice qué componente se ha
analizado, cuántos errores tiene y con qué perfil. Después vienen los errores
por archivo y por regla, el identificador de la tarea y, siempre en la última
línea, \textbf{\code{Next step:}}, lo que el asistente hará a continuación. Ese
paso se calcula a partir del estado del código, así que no tienes que decidirlo
tú. Las respuestas están en inglés porque son instrucciones para el modelo; el
asistente te las resume en tu idioma.
\end{note}

\subsubsection{3. Las reparaciones automáticas}
El asistente sigue el paso y llama a \tool{sheriff_fix}. Sheriff ordena el
import por sí solo y lista lo que queda, con el código que acepta para cada
regla:

\begin{answer}
Ran every fixer Sheriff has for these errors on basket: 4 error(s) before, 3 after (1 repaired).
Files the fixers changed (1): Basket.java, under basket/src/main/java/demo/
Still outstanding: 3 error(s) in 1 file(s). This answer lists the 3 error(s) of the first 1 file(s), and the work is those files now:
Paths below are under basket/src/main/java/demo/

Basket.java
  - Braces are not used in if statement 'if (item == null)
    return;' (lines 10-10).  [JavaBracesIfStatementsMessage2]
  - Method 'add' in file 'Basket.java' does not have JavaDoc comment  [JavaJavaDocCommentInMethodMessage6]
  - Method 'size' in file 'Basket.java' does not have JavaDoc comment  [JavaJavaDocCommentInMethodMessage6]

What Sheriff accepts, for the rules above:
Every branch and loop body has braces, and its content starts on a line of its own:
    if (page < FIRST_PAGE) {
        throw new IllegalArgumentException(UNKNOWN);
    }
[...]
Task a3f3e1be.
Next step: fix by hand every error listed above, all of them, in basket/src/main/java/demo/Basket.java, each file from the last line up [...] Then call the sheriff_fix tool with component 'basket' again.
\end{answer}

\subsubsection{4. Lo que queda, a mano}
El asistente añade las llaves y los dos JavaDoc, de la última línea hacia arriba
(así no se le mueven los números de línea), y vuelve a llamar a
\tool{sheriff_fix}, que comprueba y repara lo que esas ediciones hayan podido
romper:

\begin{answer}
Ran every fixer Sheriff has for these errors on basket: 0 error(s) before, 0 after (0 repaired).

Task 50828b9f.
Next step: run the project's tests (cd '/home/ana/proyectos/basket' && mvn test). If they pass, the component is done. If you change any file to make them pass, call the sheriff_test tool with component 'basket' again.
\end{answer}

\subsubsection{5. Los tests, y terminado}
El asistente ejecuta los tests del proyecto. Si pasan, el componente está
terminado: \textbf{0 errores y tests en verde}. Te lo cuenta en esos términos y
te dice qué errores ya estaban antes de empezar.

\subsection{El ciclo de trabajo}

Cualquier petición que cambie código sigue el mismo ciclo. No tienes que
pedirlo: el servidor se lo indica al asistente al conectarse, y cada respuesta
le da el paso siguiente.

\begin{center}
\begin{tikzpicture}[node distance=14mm and 6mm]
  \node[pbstrong,minimum width=25mm,minimum height=14mm] (ask) {Tu petición};
  \node[pbbox,right=of ask,minimum width=31mm,minimum height=14mm] (test) {\textbf{1 · Comprobar}\\[1pt]\scriptsize\code{sheriff_test}};
  \node[pbbox,right=of test,minimum width=36mm,minimum height=14mm] (clean) {\textbf{2 · Limpiar lo que había}\\[1pt]\scriptsize\code{sheriff_fix} y a mano};
  \node[pbbox,below=of clean,minimum width=36mm,minimum height=14mm] (write) {\textbf{3 · Escribir el cambio}\\[1pt]\scriptsize con \code{sheriff_guidelines}};
  \node[pbbox,left=of write,minimum width=31mm,minimum height=14mm] (again) {\textbf{4 · Comprobar otra vez}\\[1pt]\scriptsize hasta 0 errores};
  \node[pbbox,left=of again,minimum width=25mm,minimum height=14mm] (tests) {\textbf{5 · Tests}\\[1pt]\scriptsize del proyecto};
  \node[pbgold,below=9mm of tests,minimum width=25mm] (done) {\textbf{Terminado}\\[1pt]\scriptsize 0 errores, tests en verde};
  \draw[pbarrow] (ask) -- (test);
  \draw[pbarrow] (test) -- (clean);
  \draw[pbarrow] (clean) -- node[pblabel,right]{0 errores} (write);
  \draw[pbarrow] (write) -- (again);
  \draw[pbarrow] (again) -- (tests);
  \draw[pbarrow] (tests) -- (done);
  \path (clean) edge[pbarrow,draw=pbaccent,loop above,min distance=9mm,in=60,out=120] node[pblabel,above]{hasta 0 errores} (clean);
  \draw[pbarrow,draw=pbaccent] ([xshift=-4mm]tests.south east) to[out=-60,in=-120,looseness=1.1] node[pblabel,below]{si cambias algo, otra vez} (again.south);
\end{tikzpicture}
\end{center}

% ---------------------------------------------------------------------------
\clearpage
\section{Jugadas del día a día}
\label{sec:jugadas}

\begin{play}{Añadir una funcionalidad o corregir un error}
\when Cualquier cambio de código en un componente que Sheriff analiza.
\asks Lo que necesites, como siempre: \ask{añade a Basket un método que
devuelva el número de artículos distintos}.
\happens
\begin{enumerate}
  \item El asistente comprueba el componente y, si tiene errores, los limpia
        primero (sección~\ref{sec:primera}).
  \item Consulta las reglas que afectan a lo que va a escribir y escribe el
        cambio.
  \item Comprueba otra vez y repara lo que el cambio haya roto.
  \item Pasa los tests del proyecto.
  \item En Java, si ha añadido un método no privado que ningún test llama,
        escribe ese test, con casos límite (una lista vacía, un \code{null}).
\end{enumerate}
\done Te informa de 0 errores y tests en verde, y de qué errores estaban antes
de empezar.
\end{play}

\begin{play}{Revisar un proyecto sin cambiar nada}
\when Quieres saber cómo está un proyecto antes de ponerte con él.
\asks \ask{revisa los estándares de código, sin cambiar nada}.
\happens \tool{sheriff_test} analiza cada componente y da el recuento por
archivo y por regla. Como has pedido expresamente que no cambie nada, el
asistente se queda ahí.
\done Tienes el recuento y el código sigue igual (\code{git status} limpio).
\end{play}

\begin{play}{Arreglar lo que Sheriff arregla solo}
\when Un componente con muchos errores de formato: imports, llaves, nombres.
Es gratis y rápido, así que conviene hacerlo antes que nada.
\asks \ask{arregla lo que Sheriff pueda arreglar solo}.
\happens \tool{sheriff_fix} aplica todas las reparaciones automáticas de
Sheriff, en rondas mientras baja el recuento, e informa de los errores antes y
después y de qué archivos ha cambiado. No usa ningún modelo, así que
\textbf{no gasta tokens}. Edita los archivos en su sitio, sin hacer commit.
\done La respuesta dice \code{N error(s) before, M after}. Revisa el cambio
con \code{git diff}; para deshacerlo, \code{git restore .}
\end{play}

\begin{play}{Consultar las reglas antes de escribir}
\when Vas a escribir algo y quieres saber qué exige Sheriff, o un error no se
entiende.
\asks \ask{¿qué reglas hay sobre javadoc?} o \ask{explícame la regla
SortedImport}.
\happens \tool{sheriff_guidelines} busca en el catálogo de reglas. Con una
palabra, las reglas que coinciden; con el código de una regla, la regla entera y
el código que Sheriff acepta para ella. La primera consulta en una máquina
tarda unos segundos más: es cuando se extrae el catálogo de la imagen.
\done Tienes la regla, su explicación y un ejemplo correcto.
\end{play}

\begin{play}{Un módulo concreto de un proyecto con varios}
\when El repositorio tiene varios módulos (\code{api}, \code{web},
\code{core}\ldots) y te interesa uno.
\asks \ask{revisa los estándares del módulo api}.
\happens El componente es siempre el nombre de la carpeta del módulo, nunca una
ruta dentro de él; si el asistente pasa una ruta, la herramienta le contesta con
el componente correcto. Los hooks solo miran los componentes que ha cambiado el
turno.
\done Igual que en la jugada 1, para ese módulo.
\end{play}

\begin{play}{Reparar un componente entero sin dirigir cada paso}
\when Un componente con cientos de errores, típicamente la primera vez que se
analiza un proyecto que ya existía, y prefieres que se repare solo mientras
haces otra cosa.
\asks \ask{repara el componente entero con sheriff\_autofix}.
\happens El asistente te pide permiso, porque esta herramienta \textbf{hace
commits y gasta tokens}. Responde al momento con un identificador de tarea y
trabaja en segundo plano:
\begin{enumerate}
  \item Crea una rama propia, \code{sheriff-agent/<fecha>}.
  \item Aplica las reparaciones automáticas de Sheriff, en un commit propio.
  \item Pasa al modelo los archivos con errores por lotes de cinco, con un commit
        por pasada, mientras el recuento baje.
  \item Si el modelo toca un archivo fuera de su lote, se detiene sin hacer
        commit y deja ese trabajo en el árbol para que lo revises.
  \item Al terminar, pase lo que pase, ejecuta los tests del proyecto.
\end{enumerate}
Para ver cómo va, pregunta \ask{¿cómo va la tarea?}: el asistente llama a
\tool{sheriff_task}.
\done La tarea aparece como terminada. Revisa la rama (\code{git log},
\code{git diff main...}) y fusiónala si te convence.
\end{play}

\begin{warning}
Deja la sesión abierta mientras la reparación trabaja. En Codex, cerrar la
sesión la corta (sección~\ref{sec:codex}).
\end{warning}

\begin{play}{Reparar desde la terminal, sin asistente}
\when Quieres lanzar el bucle de reparación tú mismo, por ejemplo en una
máquina sin ninguna sesión de asistente abierta.
\asks Desde la carpeta del módulo:
\begin{command}[basicstyle=\ttfamily\footnotesize]{bash}
# El bucle completo
java -jar ~/.local/share/sheriff-agent/sheriff-mcp.jar --agent
# Solo analizar, o solo las reparaciones automáticas
java -jar ~/.local/share/sheriff-agent/sheriff-mcp.jar --agent --check-only
java -jar ~/.local/share/sheriff-agent/sheriff-mcp.jar --agent --sheriff-fix
# Buscar reglas
java -jar ~/.local/share/sheriff-agent/sheriff-mcp.jar --agent --rules javadoc
\end{command}
\happens El bucle es el mismo que el de la jugada anterior: rama propia, un
commit por pasada y tests al final. Por defecto contesta Claude Code con tu
suscripción, sin clave de API. La variable \code{AI_BACKEND} elige otro:
\code{codex_cli}, \code{antigravity_cli} (sección~\ref{sec:antigravity}), \code{anthropic_api},
\code{openai_api} o \code{local} (un servidor compatible con OpenAI, como
Ollama).
\done El resumen final dice cuántos errores quedan y si los tests pasan.
\end{play}

\begin{play}{Impedir que entren errores en CI}
\when Quieres que la integración continua falle si el código no cumple.
\asks Desde la raíz del proyecto:
\begin{command}{bash}
java -jar ~/.local/share/sheriff-agent/sheriff-mcp.jar --check
\end{command}
\happens Analiza todos los componentes y termina con un código de salida:

\begin{center}
\rowcolors{2}{pbcode}{white}
\begin{tabularx}{0.9\linewidth}{@{\hspace{6pt}}>{\ttfamily\bfseries}p{12mm}L}
\headrow \thead{Salida} & \thead{Significado} \\
0 & Limpio. \\
1 & Hay errores; se listan. \\
2 & No se pudo comprobar: Docker no está disponible, o no hay nada del lenguaje
    del perfil que analizar. \\
\end{tabularx}
\end{center}

Si la imagen de Sheriff no está, la descarga antes de analizar. En GitHub
Actions, por ejemplo, sin token ni cuenta:

\begin{command}{GitHub Actions}
- uses: actions/setup-java@v5
  with: { distribution: temurin, java-version: '17' }
- run: |
    curl -fsSLO https://github.com/kaizten/sheriff-mcp-server/\
    releases/latest/download/sheriff-mcp.jar
- run: java -jar sheriff-mcp.jar --check
\end{command}
\done El paso de CI está en verde, o en rojo con la lista de errores.
\end{play}

\begin{play}{Impedir que entren errores en el build de Maven}
\when Un proyecto Maven en el que \code{mvn package} debe fallar si hay errores.
\asks Añade el plugin al \code{pom.xml}, sin configuración:
\begin{command}[escapeinside={(*}{*)}]{pom.xml}
<plugin>
  <groupId>com.kaizten</groupId>
  <artifactId>sheriff-maven-plugin</artifactId>
  <version>(*\pluginversion*)</version>
  <executions>
    <execution><goals><goal>check</goal></goals></execution>
  </executions>
</plugin>
\end{command}
\happens \code{sheriff:check} se ejecuta en la fase \code{package}, después de los
tests: \code{mvn package}, \code{mvn verify} y \code{mvn install} fallan si hay
errores. \code{mvn sheriff:fix} aplica las reparaciones automáticas y falla si
queda algo. Los dos dejan el JSON de Sheriff en \path{target/sheriff/}.
\done \code{mvn package} pasa, o falla nombrando los errores.
\end{play}

\begin{note}
El plugin de Maven no viene en la release: se instala en \path{~/.m2} al
compilar el repositorio de las herramientas (\code{mvn install} en un clon).
Para la CI, \code{--check} no necesita nada más que el jar.
\end{note}

% ---------------------------------------------------------------------------
\section{Los hooks: qué verás cuando te paren}
\label{sec:hooks}

Las herramientas MCP las usa el asistente cuando le parece. Los hooks, no: los
ejecutan Claude Code y Codex en momentos fijos, y son los que garantizan que el
trabajo no se da por terminado con errores. El instalador los activa en los dos,
para todos tus proyectos; en Codex, después de que los apruebes una vez
(sección~\ref{sec:codex}). Antigravity no tiene hooks de Sheriff.

\begin{center}
\rowcolors{2}{pbcode}{white}
\begin{tabularx}{\linewidth}{@{\hspace{6pt}}K{24mm}>{\raggedright\arraybackslash}p{35mm}L}
\headrow \thead{Hook} & \thead{Cuándo actúa} & \thead{Qué hace} \\
Compuerta & Antes de la primera edición de la sesión en un componente &
  Si el componente ya tiene errores, bloquea esa edición una vez, con la lista,
  para que se limpien primero. Las siguientes pasan, y una edición a un fichero
  que tiene errores nunca se bloquea: es el arreglo. Solo en Claude Code: Codex
  edita con parches, y el fin de turno detecta lo mismo después. \\
Fin de turno & Cuando el asistente va a terminar &
  Si lo que ha cambiado en el turno tiene errores, o métodos nuevos sin test
  (en Java), le devuelve el turno con la lista. \\
Inicio de turno & Al llegar tu mensaje &
  Anota cómo estaban los componentes, para que el fin de turno mire solo lo que
  cambia en este turno y no el trabajo anterior sin confirmar. \\
\end{tabularx}
\end{center}

Cuando un hook para al asistente, verás en la conversación un mensaje como
estos, seguido del mismo tipo de \code{Next step:} que dan las herramientas:

\begin{answer}[Compuerta]
Blocked: basket already has 4 Sheriff error(s), from before this change, and basket/src/main/java/demo/Discount.java
has none of them. They are fixed first, in every file of the component, before the
new code. Edits to the files that have them go through; make this one again once
they are fixed, or now if it is part of fixing them.
\end{answer}

\begin{answer}[Fin de turno]
Not done: what this turn changed has Sheriff errors, those that were already there
included. All of them are fixed before finishing; only a false positive of the checker
may stay, named, with the reason.
\end{answer}

\subsection{Cómo se comportan}

\begin{itemize}
  \item \textbf{Nunca te bloquean por un fallo suyo.} Sin Java o sin Docker, el
        trabajo sigue y queda una línea de aviso.
  \item \textbf{Donde no hay nada que Sheriff analice}, dejan pasar todo en
        una décima de segundo.
  \item \textbf{No entran en bucle.} El fin de turno devuelve el turno como mucho
        cinco veces seguidas, y si el asistente no consigue cambiar nada, lo deja
        terminar.
  \item \textbf{Se apartan para el agente de reparación}, que ya hace sus propias
        comprobaciones.
\end{itemize}

\subsection{Exigir también los tests al terminar}

Por defecto, el fin de turno exige 0 errores pero no ejecuta los tests, porque
pueden tardar minutos. Para exigirlos también, exporta esta variable antes de
abrir Claude Code. Con Codex no es posible: no pasa a sus hooks las variables de
tu terminal.

\begin{command}{bash}
export SHERIFF_STOP_RUNS_TESTS=1
claude
\end{command}

\subsection{Ver, quitar o limitar a un proyecto}

\begin{itemize}
  \item \textbf{Ver los hooks activos:} \code{/hooks}, dentro de Claude Code o de
        Codex.
  \item \textbf{Quitarlos:}
\begin{command}{bash}
java -jar ~/.local/share/sheriff-agent/sheriff-mcp.jar --uninstall-hooks --user
\end{command}
        En Codex, \code{--uninstall-codex} los quita junto con el servidor y el
        orden de trabajo (sección~\ref{sec:desinstalar}).
  \item \textbf{Tenerlos solo en un proyecto}, en Claude Code: instala con \code{--no-hooks} y,
        desde la raíz de ese proyecto, ejecuta:
\begin{command}{bash}
java -jar ~/.local/share/sheriff-agent/sheriff-mcp.jar --install-hooks
\end{command}
        Se escriben en su \path{.claude/settings.json}. Ese archivo nombra la
        ruta del jar: súbelo al repositorio solo si todo el equipo instala en la
        misma ruta.
\end{itemize}

% ---------------------------------------------------------------------------
\section{Configurar tu proyecto}
\label{sec:configurar}

No hace falta configurar nada. Esta sección es para cuando quieras algo distinto
de lo que las herramientas deducen solas.

\subsection{Qué se analiza}

\begin{center}
\rowcolors{2}{pbcode}{white}
\begin{tabularx}{\linewidth}{@{\hspace{6pt}}>{\bfseries}p{24mm}>{\ttfamily}p{24mm}>{\ttfamily}p{30mm}L}
\headrow \thead{Lenguaje} & \thead{Archivos} & \thead{Perfil base} & \thead{Arquitecturas} \\
Java & .java & JAVA & \code{JAVA_HEXAGONAL}, \code{JAVA_DDD} \\
TypeScript & .ts, .tsx & TYPESCRIPT & \code{TYPESCRIPT_HEXAGONAL} \\
Vue & .vue & VUEJS & \\
Perl & .pl, .pm & PERL\_FORMAT & \\
Python & .py & PYTHON & \\
\end{tabularx}
\end{center}

Un componente en otro lenguaje no se da por bueno: las herramientas dicen que no
había nada que analizar. Si el proyecto se construye con Maven, Gradle o un
script \code{test} de npm, sus tests se ejecutan; si no, se analiza igual y se
dice que no se ha ejecutado ningún test.

\subsection{Declarar la arquitectura}

Sin declarar nada, un proyecto Java se comprueba con \code{JAVA}: estilo y
documentación. Si sigue una arquitectura hexagonal o DDD, declárala una vez y
todas las herramientas (servidor, hooks, \code{--check} y plugin de Maven) la
tendrán en cuenta. En el \code{pom.xml}, que los módulos heredan:

\begin{command}{pom.xml}
<properties>
  <sheriff.profile>JAVA_HEXAGONAL</sheriff.profile>
</properties>
\end{command}

O, en un proyecto sin Maven, en un archivo \code{.sheriff.properties} junto a
sus fuentes:

\begin{command}{.sheriff.properties}
profile=TYPESCRIPT_HEXAGONAL
\end{command}

Las reglas de la arquitectura \textbf{se suman} a las del lenguaje, nunca las
sustituyen. Se pueden poner varias separadas por comas
(\code{JAVA_HEXAGONAL,JAVA_DDD}); cada una añade un análisis más.

\subsection{Todos los perfiles}

Sheriff tiene más perfiles que los de arriba. Esta es la lista completa de su
imagen del 8 de octubre de 2026, con el perfil base en negrita. Cualquiera se
declara igual que una arquitectura, y la herramienta le añade siempre el perfil
base de su lenguaje: declarar \code{JAVA_HEXAGONAL_REST} ejecuta \code{JAVA}
y \code{JAVA_HEXAGONAL_REST}.

\begin{center}
\rowcolors{2}{pbcode}{white}
\begin{tabularx}{\linewidth}{@{\hspace{6pt}}>{\bfseries}p{20mm}L}
\headrow \thead{Lenguaje} & \thead{Perfiles} \\
Java & \textbf{\code{JAVA}}, \code{JAVA_HEXAGONAL}, \code{JAVA_HEXAGONAL_DOMAIN}, \code{JAVA_HEXAGONAL_APPLICATION}, \code{JAVA_HEXAGONAL_REST}, \code{JAVA_HEXAGONAL_MONGODB}, \code{JAVA_HEXAGONAL_TIMESCALE}, \code{JAVA_DDD}, \code{JAVA_DDD_ENTITY}, \code{JAVA_DDD_VALUE_OBJECT}, \code{JAVA_DDD_ENUMERATE}, \code{JAVA_ENTITY}, \code{JAVA_VALUE_OBJECT}, \code{JAVA_ENUMERATE}, \code{JAVA_USE_CASE}, \code{JAVA_POM}, \code{JAVA_COMPILATION} \\
TypeScript & \textbf{\code{TYPESCRIPT}}, \code{TYPESCRIPT_HEXAGONAL}, \code{TYPESCRIPT_HEXAGONAL_DOMAIN}, \code{TYPESCRIPT_HEXAGONAL_APPLICATION}, \code{TYPESCRIPT_HEXAGONAL_HTTP}, \code{TYPESCRIPT_HEXAGONAL_VUEJS}, \code{TYPESCRIPT_ENTITY}, \code{TYPESCRIPT_VALUE_OBJECT}, \code{TYPESCRIPT_ENUMERATE}, \code{TYPESCRIPT_FILENAME}, \code{TYPESCRIPT_LOCALE}, \code{TYPESCRIPT_PACKAGE}, \code{TYPESCRIPT_COMPILATION} \\
Vue & \textbf{\code{VUEJS}}, \code{VUEJS_COMPONENT}, \code{VUEJS_VIEW}, \code{VUEJS_FILENAME} \\
Perl & \textbf{\code{PERL_FORMAT}} \\
Python & \textbf{\code{PYTHON}} \\
\end{tabularx}
\end{center}

Para ver la lista de la versión de Sheriff que tengas, pídele un perfil que no
existe. Después de una línea de error, la línea \code{Possible Values} más
larga es la de los perfiles:

\begin{command}{bash}
docker run --rm kaizten/sheriff:latest test -t LISTA
\end{command}

Lo que comprueba cada perfil lo define Kaizten. Las herramientas conocen las
reglas exactas de 12 de ellos, los que exporta el propio Sheriff.

Para probar un perfil en una sola petición, sin declararlo, nómbralo:
\ask{revisa el módulo api con JAVA\_DDD}. Las herramientas lo usan en esa
petición, y los hooks siguen con el perfil que declara el proyecto.

\subsection{Ninguna regla se puede desactivar}

Lo que el perfil contiene es lo que se comprueba. La única excepción es un falso
positivo conocido de Sheriff, \code{FilenamePascalCase} sobre
\code{package-info.java} y \code{module-info.java}, nombres que impone Java, que
las herramientas no muestran. Si el asistente encuentra otro falso positivo, solo
puede dejarlo diciendo cuál es y por qué.

\subsection{Compartir la configuración con el equipo}

Cada máquina tiene el servidor registrado por el instalador. En Claude Code,
si prefieres que sea el proyecto quien lo declare, sube a su raíz un
\code{.mcp.json} con la ruta del jar. Usa una cosa u otra: con las dos, Claude Code avisa de que hay dos
servidores \code{sheriff} en ámbitos distintos.

\begin{command}[basicstyle=\ttfamily\footnotesize]{.mcp.json}
{
  "mcpServers": {
    "sheriff": {
      "command": "java",
      "args": ["-jar", "/home/<tu-usuario>/.local/share/sheriff-agent/sheriff-mcp.jar"]
    }
  }
}
\end{command}

\subsection{Variables útiles}

Con Claude Code basta con exportarlas antes de abrirlo. Con Codex no llegan
desde tu terminal (sección~\ref{sec:codex}).

\begin{center}
\rowcolors{2}{pbcode}{white}
\begin{tabularx}{\linewidth}{@{\hspace{6pt}}>{\ttfamily\small}p{50mm}L}
\headrow \thead{Variable} & \thead{Para qué} \\
SHERIFF\_PULL=always & Descargar una imagen nueva de Sheriff al arrancar
  (con \code{never}, ni siquiera se pregunta a Docker Hub). \\
SHERIFF\_IMAGE & Fijar la imagen por su digest
  (\code{kaizten/sheriff@sha256:...}), para que las reglas no cambien hasta que
  tú lo decidas. \\
SHERIFF\_TIMEOUT & Segundos por análisis; 300 por defecto. \\
SHERIFF\_STOP\_RUNS\_TESTS=1 & El fin de turno exige también tests en verde. \\
SHERIFF\_STOP\_MAX\_BLOCKS & Cuántas veces seguidas devuelve el turno; 5 por
  defecto. \\
VERIFICATION\_TEST\_CMD & El comando de tests, si no es el que se deduce del
  build. \\
AI\_BACKEND & Quién contesta en el bucle de reparación de la terminal. \\
\end{tabularx}
\end{center}

% ---------------------------------------------------------------------------
\section{Codex}
\label{sec:codex}

Si \code{codex} está instalado, el instalador lo configura igual que Claude
Code, para todos los proyectos. Escribe tres cosas en \path{~/.codex/}, sin
tocar el resto de esos archivos:

\begin{itemize}
  \item en \code{config.toml}, el servidor, con el tiempo límite que necesita
        (Codex corta una llamada a los 60 segundos) y \tool{sheriff_fix}
        aprobada;
  \item en \code{AGENTS.md}, el orden de trabajo, que Codex lee en todos los
        proyectos;
  \item en \code{hooks.json}, los hooks de fin y de inicio de turno. La
        compuerta no, porque Codex edita con parches; el fin de turno detecta
        los mismos errores después de editar.
\end{itemize}

\subsection{Lo que te queda a ti}

\textbf{Aprobar los hooks una vez.} Codex pregunta antes de ejecutar un hook
nuevo: en la primera sesión verás \code{Hooks need review}, o puedes hacerlo con
\code{/hooks}. Hasta que los apruebes, el fin de turno no se ejecuta.

\begin{warning}
Sin el fin de turno, Codex puede dar el trabajo por terminado con errores. En
PetClinic, con una petición de un solo método, paró dos veces a mitad, con 133
y 120 errores pendientes; fue el fin de turno lo que le hizo seguir hasta 0.
\end{warning}

\subsection{Cuatro cosas que conviene saber}

\begin{itemize}
  \item \textbf{Abre Codex en la carpeta del proyecto.} Codex no dice dónde está
        el proyecto, así que se deduce de la carpeta en la que se ejecuta.
  \item \textbf{\code{codex exec} no escribe si no se le dice.} Para que pueda
        editar: \texttt{codex exec -s workspace-write "Añade un método size() a Basket"}.
  \item \textbf{\tool{sheriff_autofix} termina con la sesión.} Codex detiene sus
        servidores al cerrar la sesión, y con ellos el bucle de reparación, que
        puede dejar el repositorio en su rama con una pasada sin confirmar. Para
        reparaciones largas, usa la terminal:
\begin{command}[basicstyle=\ttfamily\footnotesize]{bash}
AI_BACKEND=codex_cli java -jar ~/.local/share/sheriff-agent/sheriff-mcp.jar --agent
\end{command}
  \item \textbf{Las variables de tu terminal no llegan} al servidor ni a los
        hooks. Usa las opciones del instalador (\code{--pull-always},
        \code{--fail-fast}) o escríbelas en la entrada del servidor en
        \code{config.toml}: \texttt{env = \{ SHERIFF\_TIMEOUT = "600" \}}.
\end{itemize}

% ---------------------------------------------------------------------------
\section{Antigravity}
\label{sec:antigravity}

Antigravity, de Google, puede trabajar con Sheriff de dos maneras: como
asistente, con el servidor MCP, igual que Claude Code y Codex, y como uno de
los modelos del bucle de reparación. En los dos casos se usa su línea de
comandos, \code{agy}.

\subsection{Como asistente}

Si \code{agy} está instalado, el instalador lo configura, para todos los
proyectos. Escribe dos cosas, sin tocar el resto de esos archivos:

\begin{itemize}
  \item en \path{~/.gemini/config/mcp_config.json}, el servidor, el mismo que
        registraría \code{agy mcp add};
  \item en \path{~/.gemini/antigravity-cli/settings.json}, el permiso para que
        las cuatro herramientas que ni hacen commits ni gastan tokens se
        ejecuten sin preguntar, como en Claude Code: \code{mcp(sheriff/sheriff_test)}
        y las otras tres. \tool{sheriff_autofix} pregunta siempre.
\end{itemize}

Si instalas Antigravity después, repite la línea de instalación o ejecuta:

\begin{command}[basicstyle=\ttfamily\footnotesize]{bash}
java -jar ~/.local/share/sheriff-agent/sheriff-mcp.jar --install-antigravity
\end{command}

Antigravity lee el orden de trabajo del servidor, así que sigue el mismo ciclo
que Claude Code (sección~\ref{sec:primera}). Lo que no tiene son los hooks: nada impide que un turno de Antigravity termine
con errores, depende de que el modelo siga los pasos.

\begin{note}
Probado el 7 de octubre en Linux, con \code{agy} 1.3.1 y el proyecto
\code{basket} de la sección~\ref{sec:primera}. Ante \ask{arregla los errores de
estándares de código de este proyecto}, llamó a \tool{sheriff_test} y a
\tool{sheriff_fix}, corrigió a mano lo que quedaba, volvió a comprobar, pasó
\code{mvn test} y terminó con 0 errores.
\end{note}

\subsection{En el bucle de reparación}

Con \code{agy} con sesión iniciada, desde la carpeta del módulo:

\begin{command}[basicstyle=\ttfamily\footnotesize]{bash}
AI_BACKEND=antigravity_cli java -jar ~/.local/share/sheriff-agent/sheriff-mcp.jar --agent
\end{command}

El bucle es el mismo que con Claude: una rama propia, primero las reparaciones
automáticas de Sheriff, un commit por pasada y los tests del proyecto al final.
\needspace{6\baselineskip}
Para elegir el modelo (la lista sale de \code{agy models}):

\begin{command}{bash}
export ANTIGRAVITY_COMMAND="agy --mode accept-edits --model claude-sonnet-5-5-high"
\end{command}

En el bucle, Antigravity está acotado:

\begin{itemize}
  \item \textbf{Solo edita dentro del componente.} Lo que quede fuera, lo
        rechaza el propio Antigravity.
  \item \textbf{Solo puede ejecutar dos comandos}: los tests del proyecto y
        \code{git mv}. Si intenta otro, o abrir una URL, Antigravity lo rechaza
        y la pasada termina; el bucle lo anota y sigue.
  \item \textbf{Tus ajustes de Antigravity no se tocan.} Cada pasada trabaja
        con una copia temporal de ellos, que se borra al terminar.
  \item \textbf{Cada pasada tiene 600 segundos.} La variable
        \code{SHERIFF_AGENT_ANTIGRAVITY_TIMEOUT} cambia ese límite.
  \item \textbf{En Windows no se ha probado}, y por eso allí no ejecuta ningún
        comando, tampoco los tests.
\end{itemize}

\begin{note}
Probado el 7 de octubre en Spring PetClinic: de 221 errores a 0 en cinco
pasadas, con los 81 tests y el formateador del proyecto en verde.
\end{note}

% ---------------------------------------------------------------------------
\section{Solución de problemas}
\label{sec:problemas}

\newcommand\problem[2]{\par\addvspace{0.5\baselineskip}\noindent\textbf{\color{pbbrand}#1}\par\nopagebreak\vspace{-0.45\baselineskip}\noindent\hangindent=1.2em\hangafter=0 #2\par}

\subsection{Al instalar}

Estos son los mensajes del propio instalador. En los cuatro primeros se detiene; en
los demás sigue y te avisa.

\problem{\code{No Java 17 or newer was found}}{No hay ningún Java 17 o posterior
ni en \code{JAVA_HOME}, ni en el \code{PATH}, ni en las carpetas donde se
instalan los JDK. Instala uno (por ejemplo, Temurin de
\path{https://adoptium.net}; en Windows,
\code{winget install EclipseAdoptium.Temurin.21.JDK}) y repite la línea. No
hace falta que esté en el \code{PATH}.}
\problem{\code{... could not be downloaded}}{No hay red, o hay un proxy que
\code{curl} o PowerShell no conocen. Con \code{--release=TAG} (\code{-Tag}),
no hay ninguna release con ese nombre. Repite la línea cuando haya red.}
\problem{\code{the jar downloaded does not match its checksum}}{La descarga se
cortó. Repite la línea.}
\problem{\code{The MCP server does not start: FAIL}}{El instalador arrancó el
servidor con el Java que encontró, como lo arrancará tu asistente, y no
respondió. Detrás va lo que dijo el servidor; con eso, repite la línea.}
\problem{\code{Docker is not installed. The install goes on}\ldots}{Las
herramientas quedan instaladas, pero no analizarán nada hasta que haya Docker.
Instálalo, o repite la línea con \code{--install-docker}.}
\problem{\code{Docker is not running. The install goes on}\ldots}{Arranca
Docker; el servidor descargará la imagen al arrancar. En Linux, si Docker sí
está en marcha, lo más probable es que tu usuario no pueda usarlo sin
\code{sudo}: añádete al grupo con \code{sudo usermod -aG docker $USER},
cierra la sesión del sistema y vuelve a entrar.}
\problem{\code{Sheriff's image could not be pulled now}}{No hay red, o, en
Windows, Docker Desktop está en modo de contenedores Windows
(\code{no matching manifest for windows}): cámbialo a contenedores Linux, su
modo por defecto. El servidor la descargará al arrancar.}
\problem{\code{claude}, \code{codex} o \code{agy is not on the PATH}}{Ese
asistente no estaba instalado, así que el instalador no lo ha configurado.
Cuando lo instales, repite la línea, o ejecuta el comando que te ha imprimido.}
\problem{\code{No assistant was found}}{No hay instalado ninguno de Claude
Code, Codex ni Antigravity, así que el instalador no ha escrito nada en la
configuración de ningún asistente: Sheriff queda solo como comprobación de CI
(\code{--check}) y como plugin de Maven. Instala uno y repite la línea.}

\subsection{Al abrir una sesión}

\problem{\code{Failed to connect} en \code{claude mcp list}, o \code{connection closed}}{El
jar no está donde dice el registro, o el Java registrado ya no está: al
actualizar Java en Windows, la carpeta del JDK cambia de nombre. Repite la
línea de instalación, que vuelve a buscar Java y comprueba que el servidor
arranca.}
\problem{En Windows, los hooks fallan con un error \code{non-blocking} y
\code{java: command not found}}{Los escribió un instalador anterior, que ponía
\code{java} sin ruta, y Claude Code ejecuta los hooks en Git Bash, donde Java
puede no estar en el \code{PATH}. Repite la línea de instalación: ahora escribe
la ruta completa de \code{java.exe}.}
\problem{Claude Code pide permiso cada vez que usa \tool{sheriff_fix}}{Los
permisos no se escribieron: repite la instalación. \tool{sheriff_autofix}
pregunta siempre, a propósito.}
\problem{Los hooks no hacen nada}{La sesión no ha recargado la configuración.
Reiníciala y compruébalos con \code{/hooks}.}

\subsection{Al analizar}

\problem{\code{Sheriff's image ... is being downloaded}}{Es el primer uso en
esta máquina: espera unos minutos.}
\problem{\code{The analyzer could not run}}{Docker se ha parado. Arráncalo y
vuelve a pedir la revisión.}
\problem{\code{Nothing was analyzed}}{El componente no tiene fuentes del
lenguaje del perfil. Nombra el componente (\ask{revisa el módulo api}) o revisa
el perfil que declara el proyecto.}
\problem{\code{No rule catalog is available}}{No había Docker cuando arrancó el
servidor. Arranca Docker: la siguiente consulta de reglas lo vuelve a intentar,
como mucho un minuto después.}
\problem{\code{a newer kaizten/sheriff:latest has been published}}{Hay reglas
nuevas. Repite la línea de instalación para descargarlas, o instala con
\code{--pull-always}.}

\subsection{Con Codex}

\problem{\code{timed out awaiting tools/call}}{La entrada del servidor no tiene
el tiempo límite. Repite la instalación, o ejecuta
\code{--install-codex} con el jar.}
\problem{\code{writing is blocked by read-only sandbox}}{\code{codex exec} es de
solo lectura por defecto: añade \code{-s workspace-write}.}
\problem{El fin de turno no actúa}{El hook no se ha aprobado. Abre una sesión
interactiva y apruébalo con \code{/hooks}.}

\subsection{Cómo informar de un problema}

Incluye tres cosas: la versión (\code{--version} con el jar), el identificador de
la tarea que aparece en la respuesta (\code{Task 4047f176.}) y, si puedes, el
archivo \path{~/.local/state/sheriff-mcp/tasks/<id>/task.json}, que guarda con
qué versión del jar y con qué imagen exacta de Sheriff se hizo el análisis. Se
conservan las 50 tareas más recientes.

% ---------------------------------------------------------------------------
\section{Preguntas frecuentes}
\label{sec:faq}

\newcommand\faq[2]{\par\addvspace{0.6\baselineskip}\noindent{\sffamily\bfseries\color{pbbrand}#1}\par\nopagebreak\vspace{-0.45\baselineskip}#2\par}

\faq{¿Esto gasta tokens?}{Solo \tool{sheriff_autofix} y el bucle de la
terminal, que llaman a un modelo. Analizar, consultar reglas y las reparaciones
automáticas de \tool{sheriff_fix} no usan ningún modelo. El trabajo a mano que
hace tu asistente con los errores que quedan sí es parte de tu sesión normal.}

\faq{¿Se copia algo en mi proyecto?}{No. Sheriff deja unos archivos temporales
mientras analiza y las herramientas los borran después. Lo único que cambia son
los archivos que reparan \tool{sheriff_fix} o el asistente, y eso lo ves con
\code{git diff}.}

\faq{¿Por qué arregla errores que no eran míos?}{Porque el componente se entrega
limpio entero, y dejar los errores \emph{viejos} es la forma más habitual de que
nunca se arreglen. Si en un caso concreto no lo quieres, dilo en la petición.}

\faq{¿Por qué tarda tanto la primera vez?}{La imagen de Sheriff ocupa unos
4\,GB y el catálogo de reglas se extrae de ella (unos 6 segundos). A partir de
ahí, un análisis tarda unos segundos.}

\faq{¿Me llegan las reglas nuevas solas?}{No, a propósito: unas reglas que
cambian a mitad de un trabajo confunden. Cuando hay una imagen nueva, se avisa
una vez por sesión. Para descargarla, repite la instalación; para que se
descargue siempre, instala con \code{--pull-always}.}

\faq{¿Y si Docker no está en marcha?}{No hay análisis. Las herramientas lo
dicen, el asistente sigue con el trabajo y pasa los tests, y los hooks te dejan
seguir.}

\faq{¿Funciona con otros asistentes?}{Cualquier cliente MCP puede usar el
servidor. Para uno que no lea las instrucciones MCP, \code{--instructions} con
el jar imprime el orden de trabajo, y \code{--call sheriff_test} ejecuta una
herramienta desde la terminal con la misma respuesta que da el servidor.}

% ---------------------------------------------------------------------------
\section{Actualizar y desinstalar}
\label{sec:desinstalar}

\textbf{Para actualizar}, repite la línea de instalación: descarga la última
versión y vuelve a escribir la configuración. Después, reinicia las sesiones
abiertas. Repítela también si actualizas o cambias de Java: los asistentes lo
llaman por su ruta completa, y la del Java anterior puede desaparecer.

En los comandos de desinstalación, si \code{java} no está en el \code{PATH},
pon en su lugar la ruta que el instalador mostró en su primera línea
(\code{Using Java ... at ...}).

\textbf{Para desinstalar}, en este orden, porque los primeros pasos usan el
jar. En Linux y macOS:

\begin{command}{bash}
java -jar ~/.local/share/sheriff-agent/sheriff-mcp.jar --uninstall-hooks --user
# Si usas Codex
java -jar ~/.local/share/sheriff-agent/sheriff-mcp.jar --uninstall-codex
# Si usas Antigravity
java -jar ~/.local/share/sheriff-agent/sheriff-mcp.jar --uninstall-antigravity
claude mcp remove --scope user sheriff
rm -rf ~/.local/share/sheriff-agent ~/.cache/sheriff-mcp ~/.local/state/sheriff-mcp
# Opcional: libera los 4 GB de la imagen
docker rmi kaizten/sheriff:latest
\end{command}

En Windows, en PowerShell:

\begin{command}[basicstyle=\ttfamily\footnotesize]{PowerShell}
$jar = "$HOME\.local\share\sheriff-agent\sheriff-mcp.jar"
java -jar $jar --uninstall-hooks --user
# Si usas Codex
java -jar $jar --uninstall-codex
# Si usas Antigravity
java -jar $jar --uninstall-antigravity
claude mcp remove --scope user sheriff
Remove-Item -Recurse -Force -ErrorAction SilentlyContinue "$HOME\.local\share\sheriff-agent"
Remove-Item -Recurse -Force -ErrorAction SilentlyContinue "$HOME\.cache\sheriff-mcp"
Remove-Item -Recurse -Force -ErrorAction SilentlyContinue "$HOME\.local\state\sheriff-mcp"
# Opcional: libera los 4 GB de la imagen
docker rmi kaizten/sheriff:latest
\end{command}

% ---------------------------------------------------------------------------
\clearpage
\section{Tarjeta de referencia rápida}
\label{sec:tarjeta}

\begin{minipage}[t]{0.48\linewidth}\raggedright
\cardhead{Qué pedir}
\begin{itemize}[leftmargin=1.1em]
  \item \ask{revisa los estándares de código}
  \item \ask{arregla lo que Sheriff pueda arreglar solo}
  \item \ask{¿qué reglas hay sobre javadoc?}
  \item \ask{revisa los estándares del módulo api}
  \item \ask{repara el componente entero con sheriff\_autofix}
  \item \ask{¿cómo va la tarea?}
  \item \ask{\ldots{} sin tocar los errores que ya había}
\end{itemize}

\cardhead{Las cinco herramientas}
\rowcolors{2}{pbcode}{white}
\begin{tabularx}{\linewidth}{@{\hspace{4pt}}>{\ttfamily\small}p{33mm}>{\small\raggedright\arraybackslash}X}
\headrow \thead{Herramienta} & \thead{Qué hace} \\
sheriff\_test & Analiza: recuento y paso siguiente. \\
sheriff\_fix & Reparaciones automáticas; lista lo que queda. \\
sheriff\_guidelines & Las reglas y el código aceptado. \\
sheriff\_autofix & Repara todo en una rama. \textbf{Gasta tokens.} \\
sheriff\_task & Estado de una tarea. \\
\end{tabularx}

\cardhead{Salida de \code{--check}}
\rowcolors{2}{pbcode}{white}
\begin{tabularx}{\linewidth}{@{\hspace{4pt}}>{\ttfamily\bfseries\small}p{8mm}>{\small}X}
\headrow \thead{} & \thead{Significado} \\
0 & Limpio. \\
1 & Hay errores. \\
2 & No se pudo comprobar. \\
\end{tabularx}
\end{minipage}%
\hfill
\begin{minipage}[t]{0.48\linewidth}\raggedright
\cardhead{Instalar o actualizar}
{\ttfamily\footnotesize curl -fsSL https://github.com/\textbackslash\\
kaizten/sheriff-mcp-server/\textbackslash\\
releases/latest/download/\textbackslash\\
install.sh | bash\par}
{\footnotesize En Windows, la línea de la sección~\ref{sec:instalacion}.\par}

\cardhead{Comandos}
{\ttfamily\footnotesize
\begin{tabular}{@{}l@{}}
\textcolor{pbmuted}{J=\textasciitilde/.local/share/sheriff-agent/sheriff-mcp.jar}\\[3pt]
java -jar \$J -{}-version\\
java -jar \$J -{}-check\\
java -jar \$J -{}-call sheriff\_test\\
java -jar \$J -{}-agent\\
AI\_BACKEND=antigravity\_cli java -jar \$J -{}-agent\\
java -jar \$J -{}-install-hooks [-{}-user]\\
java -jar \$J -{}-uninstall-hooks [-{}-user]\\
java -jar \$J -{}-install-codex\\
java -jar \$J -{}-install-antigravity\\
\end{tabular}\par}

\cardhead{Dónde está cada cosa}
\rowcolors{2}{pbcode}{white}
\begin{tabularx}{\linewidth}{@{\hspace{4pt}}>{\ttfamily\scriptsize}p{43mm}>{\small\raggedright\arraybackslash}X}
\headrow \thead{Ruta} & \thead{Qué} \\
\textasciitilde/.local/share/sheriff-agent & El jar. \\
\textasciitilde/.cache/sheriff-mcp & El catálogo de reglas. \\
\textasciitilde/.local/state/sheriff-mcp & Las tareas. \\
\textasciitilde/.claude/settings.json & Hooks y permisos. \\
\textasciitilde/.codex/ & Configuración de Codex. \\
\textasciitilde/.gemini/config/ & Servidor en Antigravity. \\
\end{tabularx}

\cardhead{Terminado significa}
\begin{itemize}[leftmargin=1.1em]
  \item \textbf{0 errores} en el componente,
  \item \textbf{tests en verde},
  \item y un test para cada método nuevo.
\end{itemize}
\end{minipage}

\end{document}
